% SGA 3, Expose III, Loop 1 English reconstruction.
% Scope: Proposition 4.36 through section 4.38 and the Expose III
% bibliography/end matter.
% Authority: Polo--Gille Exp3-14oct24.pdf, component-local pp. 73--78
% (combined-reader pp. 183--188; re-edition pp. 173--177; original-side
% pp. 155--159). Local p. 78 / combined p. 188 is the trailing blank verso.
% OCR is a locator only. English wording was compared with Jacob C.
% Reinhold's CC BY 4.0 Markdown lineage at commit
% e7a259f3f8608ad3edf9bf6eead3fd504dd2d23e. The Polo--Gille PDF controls
% prose, mathematics, numbering, and editorial notes. There are no genuine
% graphical diagrams in this component; its exact sequences are equations
% reconstructed natively in TeX.

\setcounter{footnote}{132}

% Source page: Exp3 local p. 73; combined p. 183; re-edition p. 173;
% original-side p. 155.
\medskip
\noindent\SGAInlineRefTarget{sga3:iii:proposition:4-36}\textbf{命题 4.36.}
\textit{在 \SGARefLink{sga3:iii:theorem:4-21}{4.21} 的假设下,
再设 \(Y_{\mathscr J}\) 在 \(S_{\mathscr J}\) 上平滑, 且 \(S_0\) 仿射.
\(X\) 中在 \(S\) 上平坦 (或平滑) 且约化为 \(Y_{\mathscr J}\) 的
子 \(S\)-群 \(Y\) 所成的集合,
模去 \(X\) 在 \(S\) 上诱导 \(X_{\mathscr J}\) 的单位截面的截面所给的共轭作用后,
或者为空, 或者是下述群下的一个主齐性集:}
\[
 H^1\!\left(
 Y_0,
 \bigl[\operatorname{Lie}(X_0/S_0)/\operatorname{Lie}(Y_0/S_0)\bigr]
 \otimes_{\mathscr O_{S_0}}\mathscr J
 \right).
\]

只需验证\SGARefLink{sga3:iii:corollary:4-29}{推论 4.29} 适用, 即
\[
 d^0:
 \operatorname{Hom}_{\mathscr O_{S_0}}
   (\omega^1_{X_0/S_0},\mathscr J)
 \longrightarrow
 \operatorname{Hom}_{\mathscr O_{S_0}}
   (\mathbf n_{Y_0/X_0},\mathscr J)
\]
为满射. 这由 \SGARefLink{sga3:iii:lemma:4-25:item:ii}{4.25(ii)} 的序列 \((+)\)
分裂这一事实得到, 因为 \(S_0\) 仿射.
\footnote{编者注: 这也由 \SGARefLink{sga3:iii:corollary:4-32}{4.32} 的证明得到.}

最后陈述一个由 \SGARefLink{sga3:iii:theorem:4-21}{4.21} 和
\SGARefLink{sga3:iii:proposition:4-36}{4.36} 共同得到的推论.
事实上, 此后我们仅以这一形式使用本节的一般结果.
\footnote{编者注: 例如, 第 IX 篇用 \SGARefLink{sga3:iii:corollary:4-37}{4.37}
证明陈述 \SGARefLink{sga3:ix:theorem:3:2:bis}{3.2 bis} 与
\SGARefLink{sga3:ix:theorem:3:6:bis}{3.6 bis}.}

% Source page: Exp3 local p. 74; combined p. 184; re-edition p. 174;
% original-side p. 156.
\medskip
\noindent\SGAInlineRefTarget{sga3:iii:corollary:4-37}\textbf{推论 4.37.}
\textit{设 \(S\) 是概形, \(S_0\) 是由幂零理想 \(\mathscr I\) 定义的闭子概形.
设 \(X\) 是在 \(S\) 上平滑的 \(S\)-群,
\(Y_0\) 是 \(X_0\) 的一个子 \(S_0\)-群, 且在 \(S_0\) 上平坦.}
\begin{enumerate}
\item[\SGAInlineRefTarget{sga3:iii:corollary:4-37:item:i}(i)] \textit{若 \(S_0\) 仿射,
\(Y_0\) 在 \(S_0\) 上平滑, 且对每个 \(n\geq 0\) 均有}
\[
 H^1\!\left(
 Y_0,
 \bigl[\operatorname{Lie}(X_0/S_0)/\operatorname{Lie}(Y_0/S_0)\bigr]
 \otimes_{\mathscr O_{S_0}}
 \mathscr I^{n+1}/\mathscr I^{n+2}
 \right)=0
\]
\textit{则 \(X\) 中任意两个在 \(S\) 上平坦 (或平滑) 且约化为 \(Y_0\)
的子 \(S\)-群, 都由 \(X\) 在 \(S\) 上的一个诱导 \(X_0\) 的单位截面的截面共轭.}

\item[\SGAInlineRefTarget{sga3:iii:corollary:4-37:item:ii}(ii)] \textit{若 \(Y_0\) 仿射且为有限呈示,
并且}\footnote{编者注: 根据\SGARefLink{sga3:iii:remark:4-22}{注 4.22},
我们将 \(\operatorname{Hom}_{\mathscr O_{S_0}}
(\mathbf n_{Y_0/X_0},\mathscr I^{n+1}/\mathscr I^{n+2})\) 替换为
\(\mathbf n_{Y_0/X_0}^{\vee}\otimes_{\mathscr O_{S_0}}
\mathscr I^{n+1}/\mathscr I^{n+2}\).}
\[
 H^2\!\left(
 Y_0,
 \mathbf n_{Y_0/X_0}^{\vee}
 \otimes_{\mathscr O_{S_0}}
 \mathscr I^{n+1}/\mathscr I^{n+2}
 \right)=0
\]
\textit{对每个 \(n\geq 0\) 均成立, 则存在 \(X\) 的一个子 \(S\)-群,
它在 \(S\) 上平坦且约化为 \(Y_0\).}
\end{enumerate}

\medskip
\noindent\SGAInlineRefTarget{sga3:iii:paragraph:4-38}\textbf{4.38.}
还需联系本篇所作的三个构造. 为避免无关紧要的复杂情况, 我们取以下情形:
\(S_0\) 是域 \(k\) 的谱, \(S\) 是对偶数环 \(D(k)\) 的谱,
\(G\) 是在 \(S\) 上平滑的 \(S\)-群,
\(K\) 是一个在 \(S\) 上平滑且仿射的子 \(S\)-群.
\footnote{编者注: 下文对原文作了少量修改.
特别地, 将 \(X\) 换成 \(G\), 将 \(Y\) 换成 \(K\),
并将它们的 Lie 代数分别记为 \(\mathfrak g_0\) 和 \(\mathfrak k_0\).
还将原文使用的缩写 \(H^i({\cdot})\) 明确写为 \(H^i(K_0,{\cdot})\).}

置
\[
 \mathfrak g_0=\operatorname{Lie}G_0
 =\Gamma(S_0,\operatorname{Lie}G_0)
 =\operatorname{Lie}(G_0/S_0)(S_0),
 \qquad
 \mathfrak k_0=\operatorname{Lie}K_0.
\]
有 \(k\)-向量空间正合序列
\footnote{编者注: 配备 \(K_0\) 的伴随作用.}
\[
 0\longrightarrow\mathfrak k_0
 \xrightarrow{\ i\ }\mathfrak g_0
 \xrightarrow{\ d\ }\mathbf n^{\vee}_{K_0/G_0}
 \longrightarrow 0,
\]
它诱导出上同调正合序列
\[
\begin{aligned}
0&\longrightarrow H^0(K_0,\mathfrak k_0)
 \xrightarrow{\ i^0\ }H^0(K_0,\mathfrak g_0)
 \xrightarrow{\ d^0\ }H^0(K_0,\mathfrak g_0/\mathfrak k_0)\\
&\xrightarrow{\ \partial^0\ }H^1(K_0,\mathfrak k_0)
 \xrightarrow{\ i^1\ }H^1(K_0,\mathfrak g_0)
 \xrightarrow{\ d^1\ }H^1(K_0,\mathbf n^{\vee}_{K_0/G_0})
 \xrightarrow{\ \partial^1\ }H^2(K_0,\mathfrak k_0).
\end{aligned}
\]

这些群都有几何意义.

\begin{enumerate}
\item[\textbf{a)}]
由 II, \SGARefLink{sga3:ii:theorem:5-2-3}{5.2.3},
\(H^0(K_0,\mathfrak k_0)
=\operatorname{Lie}\underline{\operatorname{Centr}}(K_0)\).
\footnote{编者注: 由于中心化子的形成以及正规化子的形成都与基变换可交换
(参见 I, \SGARefLink{sga3:i:scholium:2-3-3-1}{2.3.3.1}),
这里用 \(\underline{\operatorname{Centr}}(K_0)\) 替换了原文中的
\(\underline{\operatorname{Centr}}(K)_0\),
同样也用 \(\underline{\operatorname{Centr}}_{G_0}(K_0)\) 和
\(\underline{\operatorname{Norm}}_{G_0}(K_0)\) 替换了
\(\underline{\operatorname{Centr}}_G(K)_0\) 和
\(\underline{\operatorname{Norm}}_G(K)_0\).}

% Source page: Exp3 local p. 75; combined p. 185; re-edition p. 175;
% original-side pp. 157--158.
\item[\textbf{b)}]
\(H^0(K_0,\mathfrak g_0)
=\operatorname{Lie}\underline{\operatorname{Centr}}_{G_0}(K_0)\)
\footnotemark[138] (同上).

\item[\textbf{c)}]
\(H^0(K_0,\mathfrak g_0/\mathfrak k_0)
=\operatorname{Lie}\underline{\operatorname{Norm}}_{G_0}(K_0)
/\mathfrak k_0\)\footnotemark[138] (同上).

\item[\textbf{d)}]
\[
 H^1(K_0,\mathfrak k_0)
 =\operatorname{Lie}\!\left(
   \underline{\operatorname{Aut}}_{S_0\text{-gr.}}(K_0)
  \right)/\operatorname{Im}(\mathfrak k_0),
\]
其中 \(\operatorname{Im}(\mathfrak k_0)\) 表示 \(\mathfrak k_0\)
在由
\[
 \operatorname{Int}_0:K_0\longrightarrow
 \underline{\operatorname{Aut}}_{S_0\text{-gr.}}(K_0).
\]
导出的态射 \(\operatorname{Lie}(\operatorname{Int}_0)\) 下的像.
事实上, 将 \SGARefLink{sga3:iii:theorem:2-1:item:ii}{2.1(ii)} 用于 \(Y=X=K\)
及 \(f_0=\operatorname{id}_{K_0}\), 可知 \(Z^1(K_0,\mathfrak k_0)\)
是 \(S_0\)-群 \(K_0\) 的无穷小自同构组成的群,
而 \(H^1(K_0,\mathfrak k_0)\) 由此群对内无穷小自同构取商得到,
即对 \(\mathfrak k_0\) 的像取商得到. 此外, 由 II,
\SGARefLink{sga3:ii:corollary:4-2-2}{4.2.2}, 还有
\[
 Z^1(K_0,\mathfrak k_0)
 =\operatorname{Lie}\!\left(
   \underline{\operatorname{Aut}}_{S_0\text{-gr.}}(K_0)
  \right).
\]
\footnote{编者注: 这是 \(\mathfrak k_0\) 的导子组成的 Lie 代数
\(\operatorname{Der}_k(\mathfrak k_0)\);
因此 \(H^1(K_0,\mathfrak k_0)\) 是
\(\operatorname{Der}_k(\mathfrak k_0)\) 对内导子取商所得的商,
也就是对 \(\operatorname{ad}:\mathfrak k_0\to\operatorname{Der}_k(\mathfrak k_0)\)
的像取商.}

\item[\textbf{e)}]
由 \SGARefLink{sga3:iii:theorem:2-1:item:ii}{2.1(ii)},
\(H^1(K_0,\mathfrak g_0)\) 是延拓典范浸入 \(i_0:K_0\to G_0\)
的同态 \(K\to G\) 之间的偏差组成的群,
模去由 \(G\) 的内自同构作用所得的偏差;
这里的内自同构由 \(G(S)\) 中在 \(G(S_0)\) 中给出单位元的元素定义
(也就是由 \(\mathfrak g_0\) 的元素定义).

\item[\textbf{f)}]
由 4.36, \(H^1(K_0,\mathbf n^{\vee}_{K_0/G_0})\) 是 \(G\) 中延拓 \(K_0\)
且在 \(S\) 上平坦 (因而在 \(S\) 上平滑, 参见 SGA 1, II, 4.10)
的子群 \(K'\) 之间的偏差组成的群,
模去由上述构造的 \(G\) 的内自同构作用所得的偏差.

\item[\textbf{g)}]
由 3.5(ii), \(H^2(K_0,\mathfrak k_0)\) 是定义在 \(K\) 上、延拓 \(K_0\) 上群结构的
各种群结构之间的偏差组成的群,
模去 \(K\) 上在 \(K_0\) 上诱导恒等映射的 \(S\)-自同构.
\end{enumerate}

现在我们要借助刚才给出的几何解释, 明确描述前面正合序列中的六个态射.

\begin{enumerate}
\item[\textbf{1)}]
态射 \(i^0\) 与 \(d^0\) 就是对典范单态射
\[
 \underline{\operatorname{Centr}}(K_0)
 \longrightarrow
 \underline{\operatorname{Centr}}_{G_0}(K_0)
 \longrightarrow
 \underline{\operatorname{Norm}}_{G_0}(K_0).
\]
取 Lie 代数 (对于 \(d^0\), 随后还要取商) 而得到的态射.
这直接由 \textbf{a)}, \textbf{b)} 与 \textbf{c)} 中所作等同的定义得到.

\item[\textbf{2)}]
态射 \(\partial^0\) 按如下方式构造. 设
\[
 \bar x\in
 \operatorname{Lie}\underline{\operatorname{Norm}}_{G_0}(K_0)
 /\mathfrak k_0.
\]
将它提升为元素
\[
 x\in\operatorname{Lie}\underline{\operatorname{Norm}}_{G_0}(K_0)
 \subset\underline{\operatorname{Norm}}_G(K)(S).
\]
则 \(\operatorname{Int}(x)\) 定义 \(K\) 的一个在 \(K_0\) 上诱导恒等映射的自同构,
因而定义 \(\operatorname{Lie}\underline{\operatorname{Aut}}_{S_0\text{-gr.}}(K_0)\)
的一个元素. 以 \(\overline{\operatorname{Int}(x)}\) 表示该元素在
\[
 \operatorname{Lie}\underline{\operatorname{Aut}}_{S_0\text{-gr.}}(K_0)
 /\operatorname{Im}(\mathfrak k_0).
\]
中的像. 则
\begin{equation}
 \partial^0(\bar x)
 =-\overline{\operatorname{Int}(x)}
 =\overline{\operatorname{Int}(x^{-1})}.
 \tag{*}\label{sga3:iii:paragraph:4-38:equation:star}
\end{equation}

事实上, 来计算由 \(\operatorname{Int}(x)\) 定义的
\(\operatorname{Lie}\underline{\operatorname{Aut}}_{S_0\text{-gr.}}(K_0)\) 的元素.
由定义, 它对应于一个元素 \(a\in Z^1(K_0,\mathfrak k_0)\), 使得
\[
 xyx^{-1}=a(y_0)y,
 \qquad y\in K(S'),\quad S'\longrightarrow S.
\]
但这也可写成
\[
 a(y_0)=xyx^{-1}y^{-1}=x-\operatorname{Ad}(y)x
 =-\partial(x)(y_0),
\]
因而 \(a=-\partial(x)\).

% Source page: Exp3 local p. 76; combined p. 186; re-edition p. 176;
% original-side p. 159.
另一方面,\footnote{编者注: 以下这句话为增补内容.} 元素
\[
 x\in\operatorname{Lie}\underline{\operatorname{Norm}}_{G_0}(K_0)
 \subset\mathfrak g_0
\]
在 \(\partial\) 下的像是 \(Z^1(K_0,\mathfrak k_0)\) 的一个元素,
其在 \(H^1(K_0,\mathfrak k_0)\) 中的像 \(\overline{\partial(x)}\)
只依赖于 \(\bar x\). 根据连接映射 \(\partial^0\) 的定义,
\[
 \partial^0(\bar x)=\overline{\partial(x)}.
\]
结合 \(a=-\partial(x)\), 便证明了 \((*)\).

\item[\textbf{3)}]
\footnote{编者注: 下文对原文作了扩充,
并在 \((**)\) 中将 \(u\circ i\) 更正为 \(i\circ u\).}
以 \(i:K\to G\) 表示典范浸入.
设 \(\bar u\in H^1(K_0,\mathfrak k_0)\) 是元素
\[
 u\in
 \operatorname{Lie}\underline{\operatorname{Aut}}_{S_0\text{-gr.}}(K_0)
 \subset\underline{\operatorname{Aut}}_{S\text{-gr.}}(K).
\]
的像. 则
\begin{equation}
 i^1(\bar u)=\bar d(i,i\circ u),
 \tag{**}\label{sga3:iii:paragraph:4-38:equation:double-star}
\end{equation}
其中 \(\bar d(i,i\circ u)\) 是 \SGARefLink{sga3:iii:remark:2-1-1}{2.1.1} 中定义的类.

事实上, \(\bar u\) 是一个元素 \(v\in Z^1(K_0,\mathfrak k_0)\) 的像,
该元素满足 \(u(y)=v(y_0)y\), 而 \(i^1(\bar u)\) 是余循环
\(i\circ v\in Z^1(K_0,\mathfrak g_0)\) 在
\(H^1(K_0,\mathfrak g_0)\) 中的像. 由于 \(i\) 为群态射,
等式 \(u(y)=v(y_0)y\) 蕴含
\[
 i(u(y))=i(v(y_0))i(y).
\]
因此 \(i\circ v=d(i,i\circ u)\), 这就证明了 \((**)\).

\item[\textbf{4)}]\SGAInlineRefTarget{sga3:iii:paragraph:4-38:item:4}
设 \(i':K\to G\) 是提升 \(i_0\) 的群态射,
\(\bar d(i,i')\) 为 \SGARefLink{sga3:iii:remark:2-1-1}{2.1.1} 中定义的类,
并设
\[
 d\bigl(K,i'(K)\bigr)
 \in C^1(K_0,\mathbf n^{\vee}_{K_0/G_0})
\]
为 \SGARefLink{sga3:iii:remark:4-5-1}{4.5.1} 中定义的偏差.
由 \SGARefLink{sga3:iii:theorem:4-21}{4.21},
\(d(K,i'(K))\) 属于 \(Z^1(K_0,\mathbf n^{\vee}_{K_0/G_0})\).
以 \(\bar d(K,i'(K))\) 表示它在 \(H^1(K_0,\mathbf n^{\vee}_{K_0/G_0})\)
中的像. 则由 \SGARefLink{sga3:iii:lemma:4-27:item:i}{4.27(i)}, 有
\begin{equation}
 d^1\bigl(\bar d(i,i')\bigr)=\bar d\bigl(K,i'(K)\bigr).
 \tag{\(\dagger\)}\label{sga3:iii:paragraph:4-38:equation:dagger}
\end{equation}
\par\smallskip
\noindent\emph{底本说明.} 在这三个相邻的出现位置, 法文重版本均印为
\(\mathbf n_{K_0/X_0}\), 未取对偶.
英文版采用 \(\mathbf n^{\vee}_{K_0/G_0}\),
这是 \SGARefLink{sga3:iii:paragraph:4-38}{4.38} 开头建立的正合序列及上同调序列所要求的.

\item[\textbf{5)}]\SGAInlineRefTarget{sga3:iii:paragraph:4-38:item:5}
最后, 设 \(K'\) 是 \(G\) 中提升 \(K_0\) 且在 \(S\) 上平坦, 因而平滑的子群.
我们已经假设 \(K_0\) 仿射. 因此知道 \(K\) 与 \(K'\)
作为延拓 \(K_0\) 的概形同构 (参见 \SGARefLink{sga3:iii:proposition:0-15}{0.15}),
所以存在 \(S\)-概形同构
\[
 f:K\xrightarrow{\ \sim\ }K'
\]
在 \(K_0\) 上诱导恒等映射. 通过 \(f\) 转移 \(K'\) 的群结构,
以 \(K_1\) 表示所得的群 (其底概形因而为 \(K\));
也就是说, \(K_1\) 的群律 \(\mu_1\) 对每个 \(S'\to S\)
和每个 \(x,y\in K(S')\) 由下式定义:
\[
 \mu_1(x,y)=f^{-1}\bigl(f(x)f(y)\bigr).
\]
由 \SGARefLink{sga3:iii:remark:3-5-1}{3.5.1},
\footnote{编者注: 下文对原文作了修改,
以计入 \SGARefLink{sga3:iii:remark:3-5-1}{3.5.1} 与
\SGARefLink{sga3:iii:proposition:4-35-1}{4.35.1} 中增补的内容.}
\(K_1\) 定义了一个余循环 \(\delta(K,K_1)\in Z^2(K_0,\mathfrak k_0)\),
使得对每个 \(S'\to S\) 和每个 \(x,y\in K(S')\), 有
\[
 \delta(K,K_1)(x_0,y_0)xy
 =\mu_1(x,y)=f^{-1}\bigl(f(x)f(y)\bigr).
\]
则由 \SGARefLink{sga3:iii:proposition:4-35-1}{4.35.1}, 有
\begin{equation}
 \partial^1\bar d(K,K')=\bar\delta(K,K_1).
 \tag{\(\ddagger\)}\label{sga3:iii:paragraph:4-38:equation:ddagger}
\end{equation}
\end{enumerate}

% Source page: Exp3 local p. 77; combined p. 187; re-edition p. 177.
% The bibliography carries no original-side page number in the re-edition.
\section*{参考文献}
\addcontentsline{toc}{section}{参考文献}
\footnote{编者注: 本篇引用的补充参考文献.}

\begin{description}
\item[{[BAlg]}] N. Bourbaki, \emph{Alg\`ebre}, 第 I--III 章,
  Hermann, 1970.
\item[{[BAC]}] N. Bourbaki, \emph{Alg\`ebre commutative}, 第 I--IV 章,
  Masson, 1985.
\item[{[DG70]}] M. Demazure 和 P. Gabriel,
  \emph{Groupes alg\'ebriques}, Masson \& North-Holland, 1970.
\item[{[Fr64]}] P. Freyd, \emph{Abelian categories}, Harper and Row,
  1964.
\end{description}

% Exp3 local p. 78 / combined p. 188 is blank.
% End of Expose III.
